\documentclass[addpoints,12pt]{exam}

\usepackage{geometry}
 \geometry{
 a4paper,
 total={170mm,260mm},
 left=20mm,
 top=20mm,
 }

\usepackage{amsmath,amssymb,amsthm}
\DeclareFontEncoding{T1}
\fontencoding{T1}
\usepackage[T2A]{fontenc} 
\usepackage[utf8]{inputenc}
\usepackage[russian]{babel}
\usepackage{graphics,graphicx}
\usepackage{indentfirst}

\usepackage{verbatim}
\usepackage{fancyvrb}  
\usepackage{booktabs}
\usepackage{multirow}

\usepackage{listings}
\usepackage{color}

\renewcommand{\solutiontitle}{\noindent\textbf{Решение:}\par\noindent}

\definecolor{gray}{gray}{0.5}

\lstset{basicstyle=\ttfamily,
    language=Matlab,    
    keepspaces=true,
    extendedchars=\true,
    identifierstyle={ \bf \color{black} },
    showstringspaces=false,
    numbers=left,%
    numbersep=7pt,  
    emph=[1]{case,switch,otherwise,nonlocal,as,yield,zip,print},    
    frame=single,    
    xleftmargin=0.3cm,
    frame=l,framesep=4pt,framerule=0.5pt,
    numberstyle={\scriptsize \color{gray}}
}


\pagestyle{headandfoot}
\runningheadrule
\firstpageheader{ИнтМатПак}{MATLAB 3}{10 марта 2026}
\runningheader{ИнтМатПак}
{MATLAB 3, Стр. \thepage\ из \numpages}
{10 марта 2026}
\firstpagefooter{}{}{}
\runningfooter{}{}{}
\pointname{} 
\boxedpoints
%\pointsinmargin
\pointsinrightmargin


%\header{MATLAB 1}
%}
%{
%	\hfill 29.05.2023
%}

% \footer{Заведующий кафедрой математических методов в экономике}
% {}
% {
%     \hfill проф. Гераськин М. И.
% }


\begin{document}
%\begin{center}
%\fbox{\fbox{\parbox{5.5in}{\centering
%Правильный.}}}
%\end{center}
\vspace{0.1in}
\makebox[\linewidth]{Ф. И. О. \enspace \hrulefill \enspace  группа: \enspace\hrulefill}
\vspace{0.1in}

\begin{questions}

\question
Поведение динамической системы описывается дифференциальным уравнением второго порядка
$$
    m \ddot{x} + a x = 2 x \sin(bt)
$$
\begin{parts}  

\renewcommand\partlabel{а)}    
\part[2]
Разработайте программу для численного интегрирования этого дифференциального уравнения. Напишите файл-функцию правых частей: 
\fillwithdottedlines{8cm}


\renewcommand\partlabel{б)}
\part[1]
Постройте график $\dot{x}(t)$ и напишите код, который определяет максимум этой функции на интервале от 3 до 5 c для $a = 5$, $b = 8$, $m = 4$ и начальных условий $x(0)=1$, $\dot{x}(0)=0$:
\fillwithdottedlines{4cm}

\renewcommand\partlabel{в)}
\part[1]
Чему равно значение локального максимума функции $\dot{x}(t)$ на интервале от 3 до 5 с? Ответ запишите с точностью до 4 значащих цифр.
\answerline[1.126]

\renewcommand\partlabel{г)}
\part[1]
Постройте фазовый портрет $(x,\dot{x})$ для $t \in [0,10]$. Чему равно расстояние от точки $[x(t_1),\dot{x}(t_1)]$ до начала координат фазовой плоскости при $t_1=10$ c? Ответ запишите с точностью до 3 значащих цифр.
\answerline[1.11]

\end{parts}

\end{questions}

\hqword{Вопрос:}
\hpword{Баллы:}
\hsword{Ответ:}
\htword{Сумма}

\vfill

\begin{center}
\gradetable[h][questions]
\end{center}

\end{document}
